% from ega4-11-fr.tex
\oldpage[IV-3]{120}

Comme les $L'_i$ sont des $A'$-modules libres, on sait (M, V, 1.1) qu'il y a
des $A'$-homomorphismes $u_i$ formant un diagramme commutatif
\[
\xymatrix{
\cdots\ar[r]&L'_{i+1}\ar[r]^{f'_i}\ar[d]_{u_{i+1}}&
L'_i\ar[r]\ar[d]_{u_i}&\cdots\ar[r]&
L'_0\ar[r]^{\varepsilon'}\ar[d]_{u_0}&M'\ar[r]\ar[d]_{1_{M'}}&0\\
\cdots\ar[r]&L''_{i+1}\ar[r]_{f''_i}&L''_i\ar[r]&\cdots\ar[r]&
L''_0\ar[r]_{\varepsilon''}&M'\ar[r]&0.
}
\]
Si l'on compose l'homomorphisme $u_\bullet:L'_\bullet\to L''_\bullet$ de
complexes ainsi défini avec l'homomorphisme canonique
$L_\bullet\to L'_\bullet$, on obtient un homomorphisme de complexes de
$A$-modules $L_\bullet\to L''_\bullet$~; en remarquant que l'on a
$(L_i\otimes_AN)\otimes_AA'=L'_i\otimes_{A'}N'$, on en déduit un
homomorphisme de complexes de $A$-modules
$L_\bullet\otimes_AN\to L''_\bullet\otimes_{A'}N'$, d'où, en passant à
l'homologie, les homomorphismes canoniques (11.2.1.1). Comme les $u_i$ sont
bien déterminés à homotopie près (M, V, 1.1), les homomorphismes (11.2.1.1)
ne dépendent pas du choix des $u_i$ ni du choix des résolutions libres
$L_\bullet$ et $L''_\bullet$.

\begin{sloppypar}
Comme les $\operatorname{Tor}_i^{A'}(M',N')$ sont des $A'$-modules, on déduit
canoniquement de (11.2.1.1) des $A'$-homomorphismes
\begin{equation}
\begin{gathered}
\varphi_i:\operatorname{Tor}_i^A(M,N)\otimes_AA'\\[-2pt]
\longrightarrow\operatorname{Tor}_i^{A'}(M',N').
\end{gathered}
\label{IV.11.2.1.5-fr}
\tag{11.2.1.5}
\end{equation}
\end{sloppypar}

Considérons maintenant deux homomorphismes d'anneaux
\[
\rho:A\to A^{(1)},\qquad \sigma:A^{(1)}\to A^{(2)},
\]
et leur composé $\sigma\circ\rho:A\to A^{(2)}$~; posons
$M^{(j)}=M\otimes_AA^{(j)}$, $N^{(j)}=N\otimes_AA^{(j)}$ pour $j=1,2$.
Alors l'homomorphisme canonique composé
\[
\operatorname{Tor}_i^A(M,N)\to
\operatorname{Tor}_i^{A^{(1)}}(M^{(1)},N^{(1)})\to
\operatorname{Tor}_i^{A^{(2)}}(M^{(2)},N^{(2)})
\]
est le même que l'homomorphisme canonique déduit de $\sigma\circ\rho$~;
cela résulte de ce que, si $L_\bullet^{(j)}$ est une résolution libre de
$M^{(j)}$, le diagramme
\[
\xymatrix{
L_\bullet\ar[r]&L_\bullet\otimes_AA^{(1)}\ar[r]\ar[d]&
L_\bullet^{(1)}\ar[d]\\
&L_\bullet\otimes_AA^{(2)}\ar[r]&
L_\bullet^{(1)}\otimes_{A^{(1)}}A^{(2)}\ar[r]&L_\bullet^{(2)}
}
\]
est commutatif.

\begin{env}[11.2.2]
\label{IV.11.2.2-fr}
Les notations étant celles de (11.2.1), considérons maintenant un
\emph{système inductif filtrant} de $A$-algèbres
$(A_\alpha,g_{\beta\alpha})$, et pour tout indice $\alpha$, posons
$M_\alpha=M\otimes_AA_\alpha$, $N_\alpha=N\otimes_AA_\alpha$~; il résulte
alors de (11.2.1) que pour chaque $i$,
$(\operatorname{Tor}_i^{A_\alpha}(M_\alpha,N_\alpha),
\psi_{i\beta\alpha})$, où $\psi_{i\beta\alpha}$ est l'homomorphisme
canonique (11.2.1.1) correspondant à $g_{\beta\alpha}$, est un \emph{système
inductif} de $A_\alpha$-modules. Posons
$A'=\varinjlim A_\alpha$, $M'=\varinjlim M_\alpha=M\otimes_AA'$,
$N'=\varinjlim N_\alpha=N\otimes_AA'$~; si on désigne par
$g_\alpha:A_\alpha\to A'$ l'homomorphisme canonique, on en déduit des
homomorphismes canoniques (11.2.1.1)
$\psi_{i\alpha}:\operatorname{Tor}_i^{A_\alpha}(M_\alpha,N_\alpha)
\to\operatorname{Tor}_i^{A'}(M',N')$ qui (toujours en vertu de (11.2.1))
forment un système inductif d'homomorphismes~; nous nous proposons de
compléter le résultat de (M, V, 9.4$^*$) en montrant que les
\begin{equation}
\psi_i=\varinjlim_\alpha\psi_{i\alpha}:
\varinjlim_\alpha\operatorname{Tor}_i^{A_\alpha}(M_\alpha,N_\alpha)
\to\operatorname{Tor}_i^{A'}(M',N')
\label{IV.11.2.2.1-fr}
\tag{11.2.2.1}
\end{equation}
sont des \emph{isomorphismes de $A'$-modules}. Pour cela, nous procédons
comme dans (M, V, 9.5$^*$), en associant à chaque $M_\alpha$ sa
\emph{résolution libre canonique}. Tout revient (compte tenu de l'exactitude
du foncteur $\varinjlim$) à prouver le
\end{env}

